\documentclass[11pt]{article}

\usepackage[margin=1in]{geometry}
\usepackage{amsmath,amssymb}

\title{Candidate Source Theorem for the Zeta Scattering Phase}
\author{Working note in \texttt{grh/}}
\date{2026-04-23}

\begin{document}

\maketitle

\begin{abstract}
This note records the exact source theorem one would want in order to close the
gap between the current RH draft's notation \(q=B_\zeta+S\) and the archive-
backed scattering quotient candidate
\[
\phi(s)=\frac{\Lambda(2s-1)}{\Lambda(2s)}.
\]
The note does not claim that this theorem has already been proved in the
canonical draft. Its purpose is to separate what is already safe at kernel
level from what still remains conditional at source level.
\end{abstract}

\section{The desired theorem}

The clean source theorem would read as follows.

\medskip
\noindent
Let
\[
\phi(s)=\frac{\Lambda(2s-1)}{\Lambda(2s)}
\]
and let \(I\subset\mathbb R\) be a compact interval avoiding the real
singularities contributed by the gamma factor, trivial zeros, and pole terms.
Choose a continuous phase \(\Phi\) on \(I\) such that
\[
\phi\!\left(\tfrac12+it\right)=e^{-2i\Phi(t)},
\qquad q(t):=\Phi'(t).
\]
Then on \(I\),
\[
q(t)=B_{\mathrm{aut}}(t)+\sum_\rho f_{\beta_\rho,\gamma_\rho}(t),
\]
where the sum runs over nontrivial zeros \(\rho=\beta_\rho+i\gamma_\rho\),
counted with multiplicity, and
\[
f_{\beta,\gamma}(t)=
\frac{1-\beta}{(1-\beta)^2+(2t-\gamma)^2}
+
\frac{\beta}{\beta^2+(2t-\gamma)^2}.
\]

This would justify the paper's source split
\[
q(t)=B_\zeta(t)+S(t),
\qquad
S(m)=q_\zeta(m)-B_\zeta(m),
\]
and identify \(S(m)\) as a positive strip-edge zero density.

\section{What is already safe}

The current repository already supports the following conservative layers.

\begin{enumerate}
\item The single-zero strip-edge kernel identity.
\item The manuscript-level formula
\[
S(m)=\sum_\rho f_{\beta_\rho,\gamma_\rho}(m)
\]
at the level of the displayed kernel formulas.
\item The conditional omitted-tail curvature theorem once the shell-convergence
package is stated.
\end{enumerate}

These are the layers isolated in \texttt{strip_edge_kernel_note.tex}.

\section{Archive provenance}

The repository's chat archive contains a direct derivation of the candidate
source theorem. In

\begin{quote}
\texttt{paper/chats/20260410-043629-69d87e03-a5c8-83a5-9f21-1062e8b6d064-riemann-hypothesis-insight.md:5528-5598}
\end{quote}

the proposed source quotient is written explicitly as
\[
\phi(s)=\frac{\Lambda(2s-1)}{\Lambda(2s)},
\qquad
\phi\!\left(\tfrac12+it\right)=e^{-2i\Phi(t)}.
\]

That archived derivation also records the single-zero pole/zero bookkeeping and
the factor-of-\(2\) normalization
\[
\partial_s\log\phi\!\left(\tfrac12+it\right)=2q(t).
\]

This is strong provenance, but it remains archive provenance until promoted into
the canonical draft.

\section{What is still missing}

The verifier concluded that the gap is not yet ``bookkeeping only.'' The paper
still lacks the theorem identifying its own abstract phase derivative
\(q=\Phi'\) with the specific quotient above.

The remaining conditional checklist is:

\begin{enumerate}
\item fix the quotient object \(\phi(s)=\Lambda(2s-1)/\Lambda(2s)\);
\item prove the branch/sign/factor-of-\(2\) normalization for the manuscript
phase;
\item prove the single-zero contribution identity inside the canonical draft;
\item state the multiplicity convention explicitly;
\item prove the compact-interval convergence/regularization statement;
\item identify and unify the full background term behind \(B_\zeta\),
\(B_{\Aut}\), and \(B_{\bg}\).
\end{enumerate}

Only after that checklist is closed does the source theorem become canonical.

\section{Use}

This note should be read as a bridge between:

\begin{enumerate}
\item the safe kernel theorem already isolated in \texttt{strip_edge_kernel_note.tex};
\item the broader candidate-object discussion in \texttt{scattering_candidate_note.tex};
\item the future task of promoting an explicit source theorem back into the main
RH draft.
\end{enumerate}

\end{document}
